% ---------
% --------
%!TEX TS-program = pdflatex
%!TEX encoding = UTF-8 Unicode
%!TEX spellcheck = English (Aspell)

\documentclass[11pt]{article}
\usepackage{lucy,handout,graphicx}
\usepackage{fourier-orns}
\usepackage{mathtools}
\usepackage[normalem]{ulem} % either use this (simple) or

\allowdisplaybreaks

\begin{document}

\noindent Name \underline{\hspace{12cm}}\\


\begin{center}
\LARGE\textbf{Quiz 2}%
\end{center}



\begin{enumerate}
%\parindent 1.5em \itemsep 4ex plus 0.5fil

%----------------------------------------------------------------------
 
\item (6 points) Suppose you have algorithms with the following runtimes.
(Assume these are exact running times as a function of the input size $n$, not
asymptotic running times.) How much slower do these algorithms get when you
double the input size? (You can find this by dividing the runtime on $2n$ by the
runtime on $n$). You should simplify your answer as much as possible, including evaluating
logarithms when possible.

\begin{enumerate}[(a)]
    \item $5n^3$ 
\vspace{1in}
    \item $2n \log_2 n$
\vspace{1in}
    \item $3^{n}$
\end{enumerate}
\vspace{1in}

\item (4 points) Simplify the following Big-O notation so that $f(n) = O\Big(g(n)\Big)$. The function $g(n)$ should be
        both as simple and as tight as possible. For example, if $f(n) = 3n^2 + 5n + 2$,
        then the answer should be $O(n^2)$, even though $O(n^3)$ is also correct but not as tight, and $O(n^2+2)$ is also
        correct but not fully simplified.

\begin{enumerate}[(a)]
    \item $f(n) = \log(6\cdot n^{n+1})$. What is $g(n)$?
\end{enumerate}
        

\end{enumerate}
%%%%%%%%%%%%%%%%%%%%
\end{document}
